\chapter{Capabilities, limits and continuing work}

The completed first release supports conversation, explicit memory, task management, current web search, public HTML reading, project code inspection, approved small-file changes, in-process history and offline direct commands. Its tested capability matrix is in docs/CAPABILITIES.md. Use it to distinguish implemented behavior from actual test evidence.

Not implemented: arbitrary command execution, automatic test execution, image, audio or PDF ingestion, login/browser automation, scheduled background jobs, reminders, multi-user service, GUI, vector search and unlimited cloud. These are future product choices, not hidden working features. The user selected the terminal version for this release.

PROJECT\_STATUS.md records the architecture, scope decision, test results, privacy constraints, report layout and optional future work. When continuing with another collaborator, read it first, then the four original agent-guidelines files. Keep reports bilingual, human-readable and attributed to the project owner. Update test evidence and documentation after each meaningful implementation change rather than claiming a feature from prompt wording alone.

A reasonable next step after using this version is collecting a few real example failures, with secrets removed, and fixing the smallest concrete problem. Larger-file patching or a graphical interface can be planned later. Do not add several frameworks and automation layers before validating the current workflow on your actual daily tasks.

All six contact URLs and the exact author name are defined centrally for the reports. Generated fonts, preview images, PDF files and personal state remain local and ignored by Git. Public source, teaching text and report sources can be reviewed before any manual commit. No commit, push, publication or deployment was performed during this work.
